\chapter{العصبونات تلعب التنس}
% full version: chapters/22_dishbrain.tex

أشهر تجربة بعصبونات حية على رقاقة هي \foreignlanguage{english}{DishBrain}، ”الدماغ في الطبق“. نحو مليون عصبون، مستنبتة على رقاقة فيها ستة وعشرون ألف قطب، لعبت أبسط تنس حاسوبي: كرة تطير على الشاشة، ومضرب يتحرك إلى الأعلى والأسفل.

\section*{كيف بُنيت اللعبة}

تدخل الكرة إلى الشبكة عبر ثمانية أقطاب. أيّها ينقر يخبر على أي ارتفاع الكرة. وكم مرة ينقر يخبر هل الكرة بعيدة: أربع مرات في الثانية حين تكون عند الجدار البعيد، ثم أكثر فأكثر، حتى أربعين، حين تقترب من المضرب. والمخرج منطقتان من الأقطاب: النقرات في إحداهما تحرك المضرب إلى الأعلى، وفي الأخرى إلى الأسفل. كل عشرة أجزاء من ألف من الثانية يعد الحاسوب النقرات ويحرك المضرب. وهذه الأجزاء العشرة هي نبضة ساعة حاسوب المنصة، لا المزرعة: الشبكة نفسها تعمل على نحو غير متزامن، ككل شبكة عصبية.

\pencilfig{22}{\pencilart{22}}{الكرة، والمضرب، ومليون عصبون تحركه.}

\section*{المعلم}

لا دوبامين على الإطلاق. بعد الإصابة تتلقى الشبكة إشارة قصيرة ومتوقعة تمامًا: الأقطاب الثمانية كلها معًا، عُشر ثانية. وبعد الإخفاق: أربع ثوانٍ من التحفيز الفوضوي في مواضع عشوائية وفي لحظات عشوائية، ثم توقف، ثم إرسال جديد للكرة.

\section*{ماذا حدث}

خلال جلسة من عشرين دقيقة طالت الجولات عند المزارع الحية، وقلت الكرات الفائتة. وكان ذلك على أفضل وجه مع التغذية الراجعة الكاملة. وحين حل الصمت ببساطة محل الإشارات بعد الإصابة والإخفاق، ظهر التحسن أيضًا لكنه كان أضعف؛ وبلا تغذية راجعة لم يظهر. ولعبت العصبونات البشرية أفضل قليلًا من عصبونات الفئران. أما التعلم الممتد على عدة أيام فلم يثبت أنه موثوق.

\section*{لماذا تتعلم الشبكة}

يفسر المؤلفون النتيجة هكذا: الشبكة تسعى إلى جعل أحاسيسها قابلة للتنبؤ. الإخفاق يجلب الفوضى، والإصابة تجلب النظام، فتتغير الشبكة كي يزداد النظام. لكن ثمة تفسير أكثر واقعية يتفق أيضًا مع البيانات: ”هُزّ بعد الإخفاق، ولا تلمس بعد الإصابة“. كل هزة تخلط الوصلات قليلًا خلطًا عشوائيًا. الإعدادات التي تخطئ فيها الشبكة كثيرًا تُخلط طوال الوقت، أما الناجحة فتبقى على حالها. فتلبث الشبكة من تلقاء نفسها حيث تقل الإخفاقات: لم يعلمها أحد، بل كفوا عن هزها فحسب. وفي نموذجنا البسيط ترفع قاعدة كهذه نسبة الإصابات فعلًا.

\pencilfig{130}{%
% scrambler: miss probability p over one setting of the network (valley in the middle); a miss kicks the setting
% at random, a hit leaves it alone, so the time spent at a setting is proportional to 1/p (6x more in the valley here)
\draw[pen] (0,1.2) -- (8.2,1.2);
\draw[penlite=pcRed] plot[smooth] coordinates {(0.0,2.46) (0.8,2.46) (1.6,2.46) (2.0,2.44) (2.4,2.41) (2.8,2.32) (3.2,2.15) (3.6,1.90) (4.0,1.62) (4.4,1.44) (4.8,1.44) (5.2,1.62) (5.6,1.90) (6.0,2.15) (6.4,2.32) (6.8,2.41) (7.2,2.44) (7.6,2.46) (8.0,2.46)};
\node[lbl, text=pcRed, anchor=south] at (7.55,2.55) {إخفاقات كثيرة};
\node[lbl, text=pcRed, anchor=west] at (5.3,1.45) {قليلة};
% kicked balls on the slopes
\draw[dotpen=pcBlue, wash=pcBlue] (1.3,2.68) circle (0.12);
\draw[arr=pcGray, densely dashed] (1.35,2.82) to[bend left=50] (2.35,2.72);
\draw[arr=pcGray, densely dashed] (1.22,2.82) to[bend right=50] (0.4,2.72);
\draw[dotpen=pcBlue, wash=pcBlue] (6.3,2.45) circle (0.12);
\draw[arr=pcGray, densely dashed] (6.25,2.59) to[bend right=50] (5.45,2.25);
\node[lbl, anchor=south] at (1.3,3.05) {هزة بعد الإخفاق};
% the ball left alone in the valley
\draw[dotpen=pcBlue, wash=pcBlue] (4.6,1.66) circle (0.12);
\node[lbl, anchor=south] at (4.6,1.95) {لا هز};
% where the network spends its time (proportional to 1/p)
\fill[pcGreen, opacity=0.22] (0,0) -- plot[smooth] coordinates {(0.0,0.18) (0.8,0.18) (1.6,0.18) (2.0,0.19) (2.4,0.19) (2.8,0.21) (3.2,0.24) (3.6,0.33) (4.0,0.55) (4.4,0.98) (4.8,0.98) (5.2,0.55) (5.6,0.33) (6.0,0.24) (6.4,0.21) (6.8,0.19) (7.2,0.19) (7.6,0.18) (8.0,0.18)} -- (8.0,0) -- cycle;
\draw[penlite=pcGreen] plot[smooth] coordinates {(0.0,0.18) (0.8,0.18) (1.6,0.18) (2.0,0.19) (2.4,0.19) (2.8,0.21) (3.2,0.24) (3.6,0.33) (4.0,0.55) (4.4,0.98) (4.8,0.98) (5.2,0.55) (5.6,0.33) (6.0,0.24) (6.4,0.21) (6.8,0.19) (7.2,0.19) (7.6,0.18) (8.0,0.18)};
\draw[pen] (0,0) -- (8.2,0);
\node[lbl, text=pcGreen!60!black, anchor=west] at (5.3,0.6) {أين تقضي الشبكة وقتها};
\node[lbl, anchor=north] at (4.1,-0.03) {إعدادات الشبكة};
}{الإخفاق يهز إعدادات الشبكة عشوائيًا، والإصابة تتركها وشأنها. حيث تكثر الإخفاقات لا تلبث الشبكة. وحيث تقل يكفون عن هزها، فتبقى هناك، مع أن أحدًا لم يقدها إلى هناك.}

كيف نميز بين التفسيرين؟ نحتاج إلى تجربة لم تُجرَ: أن تُعطى بعد الإخفاق إشارة لا فوضوية بل \emph{متوقعة} بالمدة والقوة نفسيهما. إذا تعلمت الشبكة مع ذلك، فالأمر لا يتعلق بقابلية التنبؤ، بل بمجرد الفرق بين ”بعد الإصابة“ و”بعد الإخفاق“. هذه التجربة لا تزال تنتظر من يجريها.

وهنا نهاية رحلتنا. آلة العشرين واطًا مجمّعة من قطع بسيطة، لكن كيف تجمعها بالضبط في فكرة، فلا نفهم ذلك حتى الآن إلا جزئيًا. ربما تكتب أنت الفصل التالي.

\fullversion{22}
